%\newpage
%\setcounter{page}{1}
\begin{center}
\begin{figure}
\centering%
\epsfig{file=HojaTitulo/logo.eps,scale=0.25}%
\end{figure}
\thispagestyle{empty} \vspace*{2.0cm} \textbf{\huge
An\'{a}lisis de Modos Normales \vspace{1cm}
en una Biomol\'{e}cula}\\[5.0cm]
\Large\textbf{John Erick Cabrera Ramirez}\\[6.0cm]
\small Universidad Nacional de Colombia\\
Facultad de Ciencias, Departamento de F\'{i}sica\\
Bogot\'{a}, Colombia\\
2017\\
\end{center}

%\newpage{\pagestyle{empty}\cleardoublepage}

\newpage
\begin{center}
\thispagestyle{empty} \vspace*{-0.5cm} \textbf{\huge An\'{a}lisis de Modos Normales \vspace{0.2cm} en una Biomol\'{e}cula}\\[3.0cm]
\Large\textbf{John Erick Cabrera Ramirez}\\[3.0cm]
\small Trabajo de grado presentado como requisito parcial para optar al t\'{\i}tulo de:\\
\textbf{F\'{i}sico}\\[2.5cm]
Directora:\\
PhD, Yuly Edith S\'{a}nchez Mendoza \\[2.0cm]
%T\'{\i}tulo (Ph.D., Doctor, Qu\'{\i}mico, etc.) y nombre del director(a)\\[2.0cm]
L\'{\i}nea de Investigaci\'{o}n:\\
Biof\'{i}sica Molecular\\
%Nombrar la l\'{\i}nea de investigaci\'{o}n en la que enmarca la tesis  o trabajo de investigaci\'{o}n\\
Grupo de Investigaci\'{o}n:\\
Biof\'{i}sica Molecular\\[2.5cm]
Universidad Nacional de Colombia\\
Facultad de Ciencias, Departamento de F\'{i}sica\\
Bogot\'{a}, Colombia\\
2017\\
\end{center}

\newpage
\thispagestyle{empty} \textbf{}\normalsize
\\\\\\%
\textbf{Dedicatoria}\\[4.0cm]

\begin{flushright}
\begin{minipage}{8cm}
    \noindent
        \small
        A mis padres y a mis formadores\\[1.0cm]\\
		Ya que sin ellos no hubiera sido posible realizar el proyecto de tesis tanto materialmente como en la elaboraci\'{o}n de ideas.\\
\end{minipage}
\end{flushright}

%\newpage{\pagestyle{empty}\cleardoublepage}

%\newpage
%\thispagestyle{empty} \textbf{}\normalsize
%\\\\\\%
%\textbf{\LARGE Agradecimientos}
%\addcontentsline{toc}{chapter}{\numberline{}Agradecimientos}\\\\
%Agradezco a la profesora Yuly Edith S\'{a}nchez Mendoza profesora del departamente de f\'{i}sica por ser %una excelente gu\'{i}a en el trabajo, a \\


\newpage
\textbf{\LARGE Resumen}
\addcontentsline{toc}{chapter}{\numberline{}Resumen}\\\\

%El estudio de los cotransportadores de az\'{u}car, en particular, transportadores de glucosa SGLT dependientes de sodio son esenciales en la producci\'{o}n de metabolismo y la energ\'{i}a celular.Los cotransportadores SGLT son miembros de la familia de portadores de soluto (SLC5) y algunos de estos transportadores de inter\'{e}s tienen una secuencia y estructura similar tridimensional similar. En este caso se examin\'{o} el  co-transportador dependiente  sodio galactosa del Vibrio parahaemolyticus (vSGLT), que media el transporte de galactosa en el citoplasma de las bacterias Vibrio parahaemolyticus. Seg\'{u}n la literatura, la cin\'{e}tica del co-transportador tiene entre 5 y 6 estados o conformaciones, pero en este caso de que la desvinculaci\'{o}n de los sustratos se estudia la conformaci\'{o}n, tambi\'{e}n conocido como modelo de liberaci\'{o}n estado de co-transportador que mira hacia dentro. se realiz\'{o} un estudio computacional para analizar los movimientos globales de un transportador vSGLT, y comparamos nuestros resultados computacionales con los que se encuentran en los anteriores informes experimentales. an\'{a}lisis de modos normales con un modelo el\'{a}stico de red (ENM) fue utilizado para explorar los cambios en los movimientos globales entre vSGLT en la presencia o ausencia de los iones que transportan (Na +, galactosa). ENM se ha demostrado que es un c\'{a}lculo \'{u}til herramienta para predecir la din\'{a}mica de las prote\'{i}nas de membrana en muchas aplicaciones. los modos normales m\'{a}s bajas generadas por la ENM proporcionar informaci\'{o}n valiosa sobre la din\'{a}mica global de las biomol\'{e}culas que son relevantes para su funci\'{o}n.\\

%El an\'{a}lisis de modos normales (NMA) junto con los modelos de grano grueso, han sido herramientas de la f\'{i}sica de pol\'{i}meros aplicadas posteriormente a ciertos sistemas biol\'{o}gicos donde se requiere un entendimiento del sistema a nivel molecular mas no at\'{o}mico conllevando a una  reducci\'{o}n del costo computacional. Inicialmente se estudia la teor\'{i}a y algunos de los modelos como los modelos de redes el\'{a}sticas ENM, que hacen parte del an\'{a}lisis de modos normales. En este caso, se usa como sistema biol\'{o}gico una s\'{o}la sola prote\'{i}na de membrana tipo simportador, denominada simportador de sodio/galactosa vSGLT, la cual est\'{a} presente en la membrana interna de la proteobacteria \textit{Vibrio parahaemolyticus}. Se analizan los movimientos globales del simportador mediante un modelo de red el\'{a}stica y a partir de la estructura cristalina del simportador. Se cambia la presencia de iones y substratos para ver los cambios en los movimientos globales.\\
El an\'{a}lisis de modos normales (NMA) junto con los modelos de grano grueso, han sido herramientas de la f\'{i}sica de pol\'{i}meros aplicadas posteriormente a sistemas biol\'{o}gicos como el ADN, las prote\'{i}nas, los complejos multiproteicos y  las m\'{a}quinas macromoleculares. En estos sistemas se requiere un entendimiento del sistema a nivel global mas no local o at\'{o}mico, conllevando a una  reducci\'{o}n del costo computacional. Un tipo de modelos que combina NMA con modelos de grano grueso son los modelos de redes el\'{a}sticas (ENM). En este caso, se usa como sistema biol\'{o}gico la prote\'{i}na de membrana denominada simportador de sodio/galactosa vSGLT, la cual est\'{a} presente en la membrana interna de la proteobacteria \textit{Vibrio parahaemolyticus}.\\

El estudio de los cotransportadores de az\'{u}car, en particular, los transportadores de glucosa SGLT dependientes de sodio es esencial en la producci\'{o}n de metabolismo y la energ\'{i}a celular. Los cotransportadores SGLT son miembros de la familia de simportadores de sodio soluto SSS y algunos de estos transportadores de inter\'{e}s tienen una secuencia y estructura tridimensional similar. Seg\'{u}n la literatura, la cin\'{e}tica del co-transportador tiene entre 5 y 6 estados o conformaciones, pero en este caso se estudia la liberaci\'{o}n de los sustratos en el estado de conformaci\'{o}n hacia adentro. Se analizan los movimientos globales del simportador mediante un modelo de red el\'{a}stica y a partir de la estructura cristalina del simportador. Se cambia la presencia de iones y substratos para ver los cambios en los movimientos globales.\\

\textbf{\small Palabras clave: Modos Normales, modelos de redes el\'{a}sticas, modelo de redes anisotr\'{o}picas, modelo de redes gaussianas, prote\'{i}na, cotransporte, simportador de sodio galactosa, vSGLT}.\\[2.0cm]
%\newpage
\textbf{\LARGE Abstract}\\\\
%The study of sugar cotransporters, in particular sodium-dependent glucose transporters SGLTs are essential in cellular metabolism and energy production. SGLTs are members of solute carrier family (SLC5) and some of our interest transporters have similar sequence and also 3-dimensional structure. In this case we examined the sodium- dependent glucose co-transporter of Vibrio par- ahaemolyticus (vSGLT) which mediates galactose transport into the cytoplasm of vibrio parahaemolyticus bacteria. According literature, kinetic of co- transporter has between 5 and 6 states or conformations, but in this case the un- binding of substrates is studied by inward-facing conformation, also known as state release model of co-transporter. We performed a computational study to analyze the global movements of a vSGLT transporter, and we compared our computational results with those found in previous experimental reports. Normal modes analysis with an Elastic Network Model (ENM) was used to explore the changes in global movements between vSGLT in the presence or absence of the ions they transport (Na+, galactose). ENM has been shown to be a useful computational tool for predicting the dynamics of membrane pro- teins in many applications. The lowest normal modes generated by the ENM provide valuable insight into the global dynamics of biomolecules that are rele- vant to their function.\\[2.0cm]

%The normal mode analysis (NMA) along with the coarse grain models have been tools of the polymer physics applied later to certain biological systems where an understanding of the system is required at molecular level rather than atomistic, leading to a computational cost reduction . Initially, we study the theory and some models, such as  ENM elastic network models, which are part of the normal mode analysis. In this case, a single symporter-type membrane protein, called sodium / galactose symporter vSGLT, which is present on the inner membrane of the \textit{Vibrio parahaemolyticus} proteobacteria, is used as the biological system. The global movements of the symporter are analyzed by an elastic network model and from the crystal structure of the symporter. The presence of ions and substrates is changed to seek changes in global movements.\\
Normal mode analysis (NMA) along with coarse grain models have been tools of polymer physics later applied to biological systems such as DNA, proteins, multiprotein complexes and macromolecular machines. These systems require an understanding of the system at a global level rather than local or at atomic level leading a reduction in computational cost. One type of model that combines NMA with coarse-grained models are ENM elastic network models. In this case, a membrane protein, called sodium/galactose symporter vSGLT, which is present on the inner membrane of the \textit{Vibrio parahaemolyticus} proteobacteria, is used as the biological system.\\

The study of sugar cotransporters, in particular sodium-dependent glucose transporters SGLTs is essential in cellular metabolism and energy production. SGLTs are members of Sodium-solute symporter family (SSS) and some of our interest transporters have similar sequence and also 3-dimensional structure. According literature, kinetic of co-transporter has between 5 and 6 states or conformations, but in this case the unbinding of substrates is studied by inward-facing conformation, also known as state release model of co-transporter. 
 The global movements of the symporter are analyzed by an elastic network model and from the crystal structure of the symporter. The presence of ions and substrates is changed to see changes in global motions. \\

\textbf{\small Keywords: Normal modes, elastic network models, anisotropic network model, Gaussian network model, protein, cotransporter, galactose sodium symporter, vSGLT}\\